\section{一致性检验与演化方程的实例化：软件开发人员的 $(\alpha,\beta,\mu)$}

在通过"权威评估 $\rightarrow$ 能力维度 $\rightarrow$ 语义匹配 $\rightarrow$ 离散层级 $(s,c)$"对任务影响进行重标记后，
我们将任务DNA聚合回职业级冲击参数 $(\alpha,\beta,\mu)$ 并首先进行一致性检验，
防止下游动力学依赖于错误实现的净冲击项。

\subsection{聚合值与一致性检验}

根据定义，
\begin{equation}
    \mu=\beta-\alpha.
\end{equation}
给定预览值 $\alpha=0.379113$ 和 $\beta=0.629112$，我们必须有
\begin{equation}
    \mu = 0.629112-0.379113 = 0.249999.
\end{equation}
然而，预览JSON报告 $\mu=0.504112$，这与模型定义相矛盾。
因此，在将 $\mu$ 用作 $E(t)$ 的输入之前，聚合输出必须明确声明 \texttt{mu\_definition} 并修正实现；
否则，演化推断将系统性地产生偏差。

\subsection{机制力量（异质性）检验：无得分坍缩}

使用加权方差
\begin{equation}
    \mathrm{Var}_w(s)=\sum_{j=1}^{N} w_j (s_j-\alpha)^2,\qquad
    \mathrm{Var}_w(c)=\sum_{j=1}^{N} w_j (c_j-\beta)^2,
\end{equation}
预览报告 $\mathrm{Var}_w(s)=\mathrm{Var}_w(c)=0.015608$，超过工程阈值（如 $0.005$）。
这表明重标记的 $(s,c)$ 保留了足够的异质性；机制可识别性未被坍缩。

\subsection{低置信度任务的最小敏感性（上界论证）}

对于低置信度匹配（如 task\_id=21666，权重 $w=0.05434$），可通过最坏情况边界论证稳健性。
由于层级为 $\{0.25,0.5,0.75\}$，每任务层级最大摆动为 $0.50$（$0.25\leftrightarrow0.75$），因此
\begin{equation}
    \Delta_{\max}(\alpha)\le w\cdot 0.50 = 0.05434\times 0.50 = 0.02717,
\end{equation}
类似地 $\Delta_{\max}(\beta)\le 0.02717$。
两个轻量级情景（删除该任务；重映射到次优维度）足以支撑可辩护的敏感性陈述。

\subsection{修正 $\mu$ 后进入职业演化}

一旦 $\mu$ 被修正，演化方程变为
\begin{equation}
    \frac{dE(t)}{dt}=E(t)\big(g + A(t)\mu\big),
\end{equation}
其中 $A(t)$ 是采纳情景（快速/中速/慢速），$g$ 是非AI基线（公共趋势估计或在相对指数约定下设为 $0$）。
相对尺度具有闭式解
\begin{equation}
    \frac{E(t)}{E(t_0)}=\exp\left(\int_{t_0}^{t}\big(g+A(\tau)\mu\big)\,d\tau\right).
\end{equation}
因此，在 $\mu$ 定义一致且 $A(t)$/$g$ 明确指定之前，我们不推进到 $E(t)$ 情景曲线，避免不可计算的叙述。
